\hypertarget{appendix-y-the-complete-guide-to-metaprogramming}{%
\section{Appendix Y: THE COMPLETE GUIDE TO
METAPROGRAMMING}\label{appendix-y-the-complete-guide-to-metaprogramming}}

If you can write a program that writes programs, you have reached
Godhood. C++ template metaprogramming is exactly that. It is a
Turing-complete functional programming language that executes entirely
during compilation.

\hypertarget{y.1-the-dark-arts-c98-template-recursion}{%
\subsection{Y.1 The Dark Arts: C++98 Template
Recursion}\label{y.1-the-dark-arts-c98-template-recursion}}

In C++98, we didn't have \passthrough{\lstinline!constexpr!} functions.
The only way to do math at compile time was to use struct inheritance
and recursive templates.

\hypertarget{the-compile-time-factorial}{%
\subsubsection{The Compile-Time
Factorial}\label{the-compile-time-factorial}}

\begin{lstlisting}[language={C++}]
// 1. The general case (recursive step)
template <int N>
struct Factorial {
    static const int value = N * Factorial<N - 1>::value;
};

// 2. The base case (stopping condition)
template <>
struct Factorial<0> {
    static const int value = 1;
};

int main() {
    // The compiler mathematically evaluates 5 * 4 * 3 * 2 * 1 
    // and literally just compiles "int x = 120;"
    int x = Factorial<5>::value; 
}
\end{lstlisting}

\textbf{Analogy}: It's like asking a nested doll a question. Doll 5 asks
Doll 4, Doll 4 asks Doll 3\ldots{} until Doll 0 answers ``1'', and the
answers bubble back up.

\hypertarget{y.2-the-renaissance-c11-type_traits}{%
\subsection{\texorpdfstring{Y.2 The Renaissance: C++11
\texttt{\textless{}type\_traits\textgreater{}}}{Y.2 The Renaissance: C++11 \textless type\_traits\textgreater{}}}\label{y.2-the-renaissance-c11-type_traits}}

C++11 gave us the \passthrough{\lstinline!<type\_traits>!} header, which
allowed us to inspect types. Instead of dealing with values
(\passthrough{\lstinline!int!}), we deal with Types
(\passthrough{\lstinline!typename!}).

\begin{lstlisting}[language={C++}]
#include <type_traits>

template <typename T>
void process() {
    if (std::is_pointer<T>::value) {
        // ...
    }
}
\end{lstlisting}

\emph{Wait!} The \passthrough{\lstinline!if!} statement above executes
at RUNTIME. Both sides of the \passthrough{\lstinline!if!} statement
must compile successfully, even if \passthrough{\lstinline!T!} is not a
pointer. This was the massive flaw of C++11 metaprogramming.

\hypertarget{y.3-the-workaround-sfinae-substitution-failure-is-not-an-error}{%
\subsection{Y.3 The Workaround: SFINAE (Substitution Failure Is Not An
Error)}\label{y.3-the-workaround-sfinae-substitution-failure-is-not-an-error}}

To fix the issue above, C++ engineers exploited a compiler rule. If the
compiler tries to instantiate a template, and the resulting code is
grammatically invalid, the compiler doesn't throw an error. It just
silently crosses that template off the list and looks for another one.

We weaponized this using \passthrough{\lstinline!std::enable\_if!}.

\begin{lstlisting}[language={C++}]
#include <type_traits>
#include <iostream>

// This template only "exists" if T is an integer
template <typename T>
typename std::enable_if<std::is_integral<T>::value>::type
print(T t) {
    std::cout << "Integer: " << t << "\n";
}

// This template only "exists" if T is a floating point
template <typename T>
typename std::enable_if<std::is_floating_point<T>::value>::type
print(T t) {
    std::cout << "Float: " << t << "\n";
}
\end{lstlisting}

\textbf{Analogy}: The ``Fake Door''. SFINAE is like drawing a door on a
wall. If you try to open it and it doesn't work, you just walk to the
next door instead of crashing the building.

\hypertarget{y.4-the-modern-elegance-c17-if-constexpr}{%
\subsection{\texorpdfstring{Y.4 The Modern Elegance: C++17
\texttt{if\ constexpr}}{Y.4 The Modern Elegance: C++17 if constexpr}}\label{y.4-the-modern-elegance-c17-if-constexpr}}

C++17 completely destroyed the need for 90\% of SFINAE tricks.
\passthrough{\lstinline!if constexpr!} evaluates at compile time. The
block that is \passthrough{\lstinline!false!} is completely ignored by
the compiler. It doesn't even check if the code inside it is valid for
type \passthrough{\lstinline!T!}.

\begin{lstlisting}[language={C++}]
template <typename T>
void print(T t) {
    if constexpr (std::is_integral_v<T>) {
        std::cout << "Integer: " << t << "\n";
    } else if constexpr (std::is_floating_point_v<T>) {
        std::cout << "Float: " << t << "\n";
    } else {
        std::cout << "Unknown type\n";
    }
}
\end{lstlisting}

Look how clean that is compared to SFINAE! It looks exactly like regular
C++ code.

\hypertarget{y.5-the-final-evolution-c20-concepts}{%
\subsection{Y.5 The Final Evolution: C++20
Concepts}\label{y.5-the-final-evolution-c20-concepts}}

\passthrough{\lstinline!if constexpr!} is great for branching
\emph{inside} a function. But what if you want to constrain the entire
class, or provide clear error messages to the user?

C++20 Concepts are the final, beautiful form of metaprogramming.

\begin{lstlisting}[language={C++}]
#include <concepts>

void print(std::integral auto t) {
    std::cout << "Integer: " << t << "\n";
}

void print(std::floating_point auto t) {
    std::cout << "Float: " << t << "\n";
}
\end{lstlisting}

That's it. It's perfectly safe, heavily optimized, and if a user tries
to pass a \passthrough{\lstinline!std::string!}, the compiler will
output a clean, 1-line error:
\passthrough{\lstinline!"Constraint not satisfied: std::string is not integral"!}.

This is the journey from C++98 to C++20. From dark, recursive hacks, to
beautiful, semantic constraints.

\begin{center}\rule{0.5\linewidth}{0.5pt}\end{center}

\begin{center}\rule{0.5\linewidth}{0.5pt}\end{center}

\hypertarget{volume-26-the-head-first-masterclass-beginner-to-godhood}{%
\section{VOLUME 26: THE ``HEAD FIRST'' MASTERCLASS (BEGINNER TO
GODHOOD)}\label{volume-26-the-head-first-masterclass-beginner-to-godhood}}

Welcome to Volume 26. In the previous 25 volumes, we covered the
technical specifications of C++ from C++98 to C++26. We covered HFT
patterns, compiler internals, and memory models.

But what if you are a beginner? What if all of that went completely over
your head?

In this volume, we hit the reset button. We are going to take the most
terrifying concepts in C++ and explain them using the \textbf{``Head
First''} method: extremely conversational, heavily reliant on real-world
analogies, and answering the ``dumb'' questions that everyone is too
afraid to ask.

We will start at absolute zero and build back up to Godhood.

\hypertarget{chapter-121-the-head-first-guide-to-memory}{%
\subsection{Chapter 121: The ``Head First'' Guide to
Memory}\label{chapter-121-the-head-first-guide-to-memory}}

\hypertarget{the-hotel-analogy}{%
\subsubsection{The Hotel Analogy}\label{the-hotel-analogy}}

Imagine your computer's RAM is a massive hotel called \textbf{The
Silicon Inn}. It has billions of rooms.

When you run a C++ program, you walk up to the front desk and say, ``I
need a room.''

\hypertarget{variables-the-rooms}{%
\paragraph{1. Variables (The Rooms)}\label{variables-the-rooms}}

\begin{lstlisting}[language={C++}]
int x = 5;
\end{lstlisting}

You just rented a standard-sized room. The hotel clerk paints a giant
``X'' on the door. Inside the room, they put a piece of paper with the
number ``5'' on it.

\hypertarget{pointers-the-room-key}{%
\paragraph{2. Pointers (The Room Key)}\label{pointers-the-room-key}}

\begin{lstlisting}[language={C++}]
int* p = &x;
\end{lstlisting}

You ask the clerk for a room key. A pointer
(\passthrough{\lstinline!p!}) is literally just a piece of plastic with
the room number engraved on it. It doesn't hold the number ``5''. It
holds the room number (e.g., Room 104).

\hypertarget{dereferencing-opening-the-door}{%
\paragraph{3. Dereferencing (Opening the
Door)}\label{dereferencing-opening-the-door}}

\begin{lstlisting}[language={C++}]
*p = 10;
\end{lstlisting}

The \passthrough{\lstinline!*!} symbol means ``Take this room key, walk
down the hallway, open the door, and change what's inside.'' You walk to
Room 104 and change the paper from ``5'' to ``10''. Now, if anyone looks
at variable \passthrough{\lstinline!x!}, they will see 10.

\hypertarget{the-stack-vs.-the-heap-the-backpack-vs.-the-storage-unit}{%
\subsubsection{The Stack vs.~The Heap (The Backpack vs.~The Storage
Unit)}\label{the-stack-vs.-the-heap-the-backpack-vs.-the-storage-unit}}

You have two places to store things at The Silicon Inn.

\textbf{The Stack (Your Backpack)} When you enter the hotel lobby (start
a function), you are wearing a backpack.

\begin{lstlisting}[language={C++}]
void my_function() {
    int local_var = 42; // Goes in the backpack
}
\end{lstlisting}

You throw \passthrough{\lstinline!local\_var!} into your backpack. It's
incredibly fast to put things in and take things out. But there's a
catch: when you leave the lobby (the function ends), a security guard
takes your backpack and throws it in the incinerator. Everything inside
is destroyed instantly and automatically.

\textbf{The Heap (The Storage Unit)} What if you buy a grand piano? It
won't fit in your backpack. What if you want to leave the piano at the
hotel for a friend who is arriving tomorrow?

You must rent a Storage Unit (The Heap).

\begin{lstlisting}[language={C++}]
void my_function() {
    int* piano = new int(42); // Rents a storage unit
}
\end{lstlisting}

You call \passthrough{\lstinline!new!}. The clerk hands you a key
(\passthrough{\lstinline!piano!}) to a permanent storage unit. You put
the piano inside. When you leave the lobby, the security guard burns
your backpack (which contains the \emph{key}), but the Storage Unit
itself is untouched.

\textbf{The Disaster (Memory Leak)}: You just lost the only key to the
storage unit, but you are still paying rent on it forever! This is a
\textbf{Memory Leak}.

\textbf{The Fix}: You must explicitly tell the clerk you are done before
you leave.

\begin{lstlisting}[language={C++}]
    delete piano; // Empties the storage unit and stops the rent
\end{lstlisting}

\begin{quote}
\textbf{Brain Power: Why not just use The Stack for everything?} The
Stack is small. Usually just 1 to 8 Megabytes. If you try to put a
10-Megabyte array into your backpack, the backpack rips open and the
program crashes immediately. This is called a \textbf{Stack Overflow}.
\end{quote}

\begin{center}\rule{0.5\linewidth}{0.5pt}\end{center}

\hypertarget{chapter-122-the-head-first-guide-to-object-oriented-c}{%
\subsection{Chapter 122: The ``Head First'' Guide to Object-Oriented
C++}\label{chapter-122-the-head-first-guide-to-object-oriented-c}}

Object-Oriented Programming (OOP) is how we build complex things without
going insane.

\hypertarget{the-factory-analogy}{%
\subsubsection{The Factory Analogy}\label{the-factory-analogy}}

Imagine you want to build cars.

\hypertarget{the-class-the-blueprint}{%
\paragraph{1. The Class (The Blueprint)}\label{the-class-the-blueprint}}

\begin{lstlisting}[language={C++}]
class Car {
private:
    Engine engine; // The messy wiring
public:
    void press_gas() { engine.inject_fuel(); }
};
\end{lstlisting}

A \passthrough{\lstinline!class!} is just a blueprint. You can't drive a
blueprint.

Notice the \passthrough{\lstinline!private!} and
\passthrough{\lstinline!public!} keywords? This is
\textbf{Encapsulation}. When you buy a car, Toyota gives you a gas pedal
(\passthrough{\lstinline!public!}). They do NOT let you manually inject
fuel into the engine cylinders (\passthrough{\lstinline!private!}). If
they did, you would explode the engine on day one. Encapsulation
protects the user from their own stupidity.

\hypertarget{the-object-the-physical-car}{%
\paragraph{2. The Object (The Physical
Car)}\label{the-object-the-physical-car}}

\begin{lstlisting}[language={C++}]
Car my_honda;
my_honda.press_gas();
\end{lstlisting}

Now you have a physical car. You can build 1,000 cars from one
blueprint.

\hypertarget{inheritance-the-specialized-blueprint}{%
\paragraph{3. Inheritance (The Specialized
Blueprint)}\label{inheritance-the-specialized-blueprint}}

You want to build a Racecar. A Racecar is exactly like a normal Car, but
it has a turbo boost. Instead of drawing a brand new blueprint from
scratch, you take the \passthrough{\lstinline!Car!} blueprint, put
tracing paper over it, and draw a turbocharger.

\begin{lstlisting}[language={C++}]
class Racecar : public Car {
public:
    void press_turbo() { ... }
};
\end{lstlisting}

\hypertarget{polymorphism-the-valet-driver}{%
\paragraph{4. Polymorphism (The Valet
Driver)}\label{polymorphism-the-valet-driver}}

This is the hardest concept for beginners. Imagine you are a Valet
Driver at a fancy hotel. Your job is simple: Drive the car into the
garage.

\begin{lstlisting}[language={C++}]
void park_car(Car* c) {
    c->press_gas();
}
\end{lstlisting}

A customer hands you the keys to a \passthrough{\lstinline!Racecar!}.
Can you park it? YES! Because a \passthrough{\lstinline!Racecar!}
\emph{is a} \passthrough{\lstinline!Car!}. It has a gas pedal. You don't
need to know how the turbo works to park it.

But what if a \passthrough{\lstinline!Racecar!}'s gas pedal works
differently than a normal \passthrough{\lstinline!Car!}'s gas pedal? In
C++, if you call \passthrough{\lstinline!c->press\_gas()!}, the compiler
will normally look at the pointer type (\passthrough{\lstinline!Car*!})
and call the normal car's gas pedal, ignoring the fact that it's
actually a racecar!

\textbf{The Fix: \passthrough{\lstinline!virtual!} functions}. By
marking \passthrough{\lstinline!virtual void press\_gas();!} in the base
class, you tell the Valet: ``Hey, before you press the gas, look inside
the glovebox. There is a sticky note (the \textbf{vtable}) that tells
you exactly which gas pedal to press for this specific vehicle.''

\begin{center}\rule{0.5\linewidth}{0.5pt}\end{center}

\hypertarget{chapter-123-the-head-first-guide-to-templates}{%
\subsection{Chapter 123: The ``Head First'' Guide to
Templates}\label{chapter-123-the-head-first-guide-to-templates}}

C++ is famous for its templates. They look scary, but they are just a
``Fill-in-the-Blanks'' form.

\hypertarget{the-cookie-cutter-analogy}{%
\subsubsection{The Cookie Cutter
Analogy}\label{the-cookie-cutter-analogy}}

Imagine you are a baker. You want to make a star-shaped cookie. You
could carve a star out of chocolate dough. Then carve a star out of
vanilla dough. Then carve a star out of strawberry dough. This is
exhausting (writing the same function for \passthrough{\lstinline!int!},
\passthrough{\lstinline!float!}, and \passthrough{\lstinline!double!}).

\textbf{The Solution:} You build a Star-Shaped Cookie Cutter (a
Template).

\begin{lstlisting}[language={C++}]
template <typename Dough>
Dough make_star(Dough d) {
    return shape_into_star(d);
}
\end{lstlisting}

When you type \passthrough{\lstinline!make\_star<int>(5)!}, the C++
Compiler literally copy-pastes your code, replaces the word
\passthrough{\lstinline!Dough!} with \passthrough{\lstinline!int!}, and
compiles a brand new function. When you type
\passthrough{\lstinline!make\_star<double>(3.14)!}, the compiler
copy-pastes it again and replaces \passthrough{\lstinline!Dough!} with
\passthrough{\lstinline!double!}.

\begin{quote}
\textbf{There are no dumb questions\ldots{}}

\textbf{Q: Doesn't that make my compiled program huge?} \textbf{A:} Yes!
This is called \textbf{Code Bloat}. If you call a template function with
50 different types, the compiler generates 50 different functions in the
final binary.

\textbf{Q: Does it slow down my program?} \textbf{A:} No! Actually, it
makes it FASTER. Because the compiler generates a specific function for
\passthrough{\lstinline!int!}, it can perfectly optimize it for
\passthrough{\lstinline!int!} at compile time. This is why C++ templates
are faster than Java Generics or Python functions.
\end{quote}

\hypertarget{c20-concepts-the-smart-cookie-cutter}{%
\subsubsection{C++20 Concepts (The Smart Cookie
Cutter)}\label{c20-concepts-the-smart-cookie-cutter}}

What happens if you try to use the Star Cookie Cutter on a bowl of Soup?
It makes a massive mess. In old C++, if you passed a
\passthrough{\lstinline!std::string!} into a math template, the compiler
would print 500 lines of horrific errors.

C++20 fixes this with \textbf{Concepts}. It adds a warning label to the
cookie cutter.

\begin{lstlisting}[language={C++}]
template <typename Dough>
requires IsSolid<Dough> // The Concept!
Dough make_star(Dough d) { ... }
\end{lstlisting}

Now, if you pass Soup, the compiler just says: ``Error: Soup is not
Solid.'' 1 line of error. Beautiful.

\begin{center}\rule{0.5\linewidth}{0.5pt}\end{center}

\hypertarget{chapter-124-the-head-first-guide-to-concurrency}{%
\subsection{Chapter 124: The ``Head First'' Guide to
Concurrency}\label{chapter-124-the-head-first-guide-to-concurrency}}

Multithreading is doing two things at once.

\hypertarget{the-restaurant-kitchen-analogy}{%
\subsubsection{The Restaurant Kitchen
Analogy}\label{the-restaurant-kitchen-analogy}}

\textbf{Single-Threaded}: You are the only chef in the kitchen. You chop
the onions, then you boil the water, then you cook the pasta. It takes
30 minutes. \textbf{Multi-Threaded}: You hire two sous-chefs. One chops
onions. One boils water. You cook the pasta. It takes 10 minutes.

\begin{lstlisting}[language={C++}]
#include <thread>

void chop_onions() { ... }
void boil_water() { ... }

int main() {
    std::thread chef1(chop_onions);
    std::thread chef2(boil_water);
    
    // The main thread waits for them to finish
    chef1.join(); 
    chef2.join();
}
\end{lstlisting}

\hypertarget{the-data-race-the-knife-fight}{%
\subsubsection{The Data Race (The Knife
Fight)}\label{the-data-race-the-knife-fight}}

What happens if Chef 1 and Chef 2 both try to grab the \emph{same} knife
at the \emph{same} millisecond? In C++, this is a \textbf{Data Race}. It
is Undefined Behavior. Your program will crash or produce garbage data.

\hypertarget{the-mutex-the-talking-stick}{%
\subsubsection{The Mutex (The Talking
Stick)}\label{the-mutex-the-talking-stick}}

To solve the knife fight, we use a \passthrough{\lstinline!std::mutex!}.
Think of it as a ``Talking Stick'' in a kindergarten class. If you are
holding the stick, you are allowed to use the knife. If someone else
wants the knife, they have to wait until you put the stick down.

\begin{lstlisting}[language={C++}]
#include <mutex>
std::mutex knife_mutex;

void chef1() {
    knife_mutex.lock();   // Grab the stick
    use_knife();
    knife_mutex.unlock(); // Put the stick down
}
\end{lstlisting}

\begin{quote}
\textbf{Godhood Tip}: NEVER call \passthrough{\lstinline!.lock()!} and
\passthrough{\lstinline!.unlock()!} manually. What if
\passthrough{\lstinline!use\_knife()!} throws an exception? The Chef
drops dead, but he is still holding the Talking Stick! The other chefs
wait forever. This is a \textbf{Deadlock}. Always use
\passthrough{\lstinline!std::lock\_guard!}. It is a robot that
automatically grabs the stick for you, and automatically returns it the
millisecond the function ends, even if the Chef dies.
\end{quote}

\begin{lstlisting}[language={C++}]
void chef1() {
    std::lock_guard<std::mutex> guard(knife_mutex); // Safe!
    use_knife();
} // Automatically unlocks here.
\end{lstlisting}

\begin{center}\rule{0.5\linewidth}{0.5pt}\end{center}

\hypertarget{chapter-125-the-head-first-guide-to-move-semantics}{%
\subsection{Chapter 125: The ``Head First'' Guide to Move
Semantics}\label{chapter-125-the-head-first-guide-to-move-semantics}}

This is the feature that makes Modern C++ fast.

\hypertarget{the-u-haul-box-analogy}{%
\subsubsection{The ``U-Haul Box''
Analogy}\label{the-u-haul-box-analogy}}

Imagine you have a giant, beautiful, intricately constructed Lego
Castle. You want to give it to your friend across the street.

\textbf{Before C++11 (The Copy Era)}: You cannot move the castle. You
must go to the store, buy 10,000 new Lego bricks, and spend 5 hours
building an \emph{exact replica} of the castle at your friend's house.
Then, you smash your original castle into pieces. This is incredibly
slow.

\textbf{After C++11 (The Move Era)}: You take the Lego Castle, put it in
a cardboard box, carry the box across the street, and hand it to your
friend. Time taken: 10 seconds.

\hypertarget{how-c-does-it}{%
\paragraph{How C++ does it}\label{how-c-does-it}}

When you pass a \passthrough{\lstinline!std::vector!} (the Lego Castle)
to a function, C++ wants to copy it by default to be safe.

If you want to ``Move'' it, you must use
\passthrough{\lstinline!std::move!}. \passthrough{\lstinline!std::move!}
is just a Shipping Label. It slaps a sticker on the vector that says:
\textbf{``I DO NOT CARE ABOUT THIS OBJECT ANYMORE. FEEL FREE TO STEAL
ITS GUTS.''}

\begin{lstlisting}[language={C++}]
std::vector<int> my_castle = {1, 2, 3, 4, 5... 1000000};

// Clones the castle. Takes 10 milliseconds.
take_castle(my_castle); 

// Slaps the shipping label on. Takes 0.0001 milliseconds.
take_castle(std::move(my_castle)); 
\end{lstlisting}

\textbf{The Aftermath}: After you move
\passthrough{\lstinline!my\_castle!}, it is an empty plot of land. It
has 0 elements. Do not try to use it again!

\begin{center}\rule{0.5\linewidth}{0.5pt}\end{center}

\hypertarget{chapter-126-the-head-first-guide-to-the-stl}{%
\subsection{Chapter 126: The ``Head First'' Guide to the
STL}\label{chapter-126-the-head-first-guide-to-the-stl}}

The Standard Template Library (STL) is your toolbox. If you try to build
a house using only a hammer (raw \passthrough{\lstinline!for!} loops and
raw arrays), it will take you a year and the house will fall down. If
you use the STL, you get nailguns, circular saws, and laser levels.

\hypertarget{the-three-components-of-the-stl}{%
\subsubsection{The Three Components of the
STL}\label{the-three-components-of-the-stl}}

\begin{enumerate}
\def\labelenumi{\arabic{enumi}.}
\tightlist
\item
  \textbf{Containers}: The tool belts that hold your data.
  (\passthrough{\lstinline!vector!}, \passthrough{\lstinline!map!},
  \passthrough{\lstinline!set!}).
\item
  \textbf{Iterators}: The measuring tapes. They allow you to point at
  specific items inside a container safely.
\item
  \textbf{Algorithms}: The power tools.
  (\passthrough{\lstinline!std::sort!},
  \passthrough{\lstinline!std::find!},
  \passthrough{\lstinline!std::reverse!}).
\end{enumerate}

\hypertarget{why-algorithms-are-better-than-for-loops}{%
\subsubsection{\texorpdfstring{Why Algorithms are better than
\texttt{for}
loops}{Why Algorithms are better than for loops}}\label{why-algorithms-are-better-than-for-loops}}

Imagine you want to find the number \passthrough{\lstinline!42!} in a
list. You could write a \passthrough{\lstinline!for!} loop. But a
\passthrough{\lstinline!for!} loop is just a loop. The person reading
your code has to read all 5 lines of the loop to figure out \emph{what}
you are trying to do.

If you use \passthrough{\lstinline!std::find(v.begin(), v.end(), 42)!},
the person reading your code instantly knows your intent. Furthermore,
\passthrough{\lstinline!std::find!} is written by C++ compiler
engineers. It is heavily optimized, unrolled, and bug-free. Your
\passthrough{\lstinline!for!} loop might have an off-by-one error.

\textbf{Godhood Tip}: The famous C++ speaker Sean Parent has a rule:
\textbf{``No Raw Loops.''} If you are writing a
\passthrough{\lstinline!for!} loop, there is almost certainly an STL
algorithm that does what you want, but safer and faster.

\begin{center}\rule{0.5\linewidth}{0.5pt}\end{center}

\begin{center}\rule{0.5\linewidth}{0.5pt}\end{center}
